\begin{tikzpicture}[font=\small,>=stealth]
\draw[gray!45,rounded corners=3pt] (0,0) rectangle (14.2,8.5);
\node at (7.1,7.85) {方法五：两焦点作圆心，交点个数决定是否有解};

\node at (7.1,6.9) {取 $a=1,\lambda=2$，于是 $r_1=4,r_2=2$；只改变 $c=e$。};
\node at (2.5,6.2) {$e=2$};\begin{scope}[shift={(2.75,4.4)},scale=.3]
\draw[orange!85!black,thick] (-2,0) circle(4);\draw[purple!80!black,thick] (2,0) circle(2);\draw[gray] (-2,0)--(2,0);\fill (-2,0) circle(3pt) node[below=43pt] {$F_1$};\fill (2,0) circle(3pt) node[below=43pt] {$F_2$};
\fill[teal!75!black] (1.5,1.9364916731037085) circle(4pt);\fill[teal!75!black] (1.5,-1.9364916731037085) circle(4pt);\end{scope}
\node at (2.5,1.65) {相交：两个点};
\node at (7.0,6.2) {$e=3$};\begin{scope}[shift={(7.25,4.4)},scale=.3]
\draw[orange!85!black,thick] (-3,0) circle(4);\draw[purple!80!black,thick] (3,0) circle(2);\draw[gray] (-3,0)--(3,0);\fill (-3,0) circle(3pt) node[below=43pt] {$F_1$};\fill (3,0) circle(3pt) node[below=43pt] {$F_2$};
\fill[teal!75!black] (1.0,0.0) circle(4pt);\end{scope}
\node at (7.0,1.65) {外切：一个顶点};
\node at (11.5,6.2) {$e=3.5$};\begin{scope}[shift={(11.75,4.4)},scale=.3]
\draw[orange!85!black,thick] (-3.5,0) circle(4);\draw[purple!80!black,thick] (3.5,0) circle(2);\draw[gray] (-3.5,0)--(3.5,0);\fill (-3.5,0) circle(3pt) node[below=43pt] {$F_1$};\fill (3.5,0) circle(3pt) node[below=43pt] {$F_2$};
\end{scope}
\node at (11.5,1.65) {分离：没有点};

\node at (7.1,.65) {三图同一尺度；圆心距 $2c$ 与半径和 $r_1+r_2=6$ 比较。};

\end{tikzpicture}
